
\فصل{تعریف مسئله، اهداف و فرصت محصول}

\قسمت{تعریف مسئله}

مدیریت مجتمع‌های مسکونی در ایران عمدتاً به‌صورت سنتی و دستی انجام می‌شود.
این رویکرد دارای چالش‌های جدی زیر است:

\شروع{فقرات}
\فقره \textbf{بی‌شفافیتی مالی:} پرداخت شارژ و پیگیری بدهی‌ها از طریق اطلاعیه‌های کاغذی یا تماس تلفنی انجام می‌شود و هیچ تاریخچه‌ی قابل اتکایی وجود ندارد.
\فقره \textbf{ناکارآمدی درخواست‌های خدماتی:} ثبت و پیگیری درخواست‌های تعمیرات به‌صورت حضوری یا تلفنی انجام می‌شود، ردیابی وضعیت آن‌ها امکان‌پذیر نیست و در صورت تغییر مدیر ساختمان، تاریخچه از بین می‌رود.
\فقره \textbf{مدیریت دستی امکانات مشترک:} رزرو استخر، سالن اجتماعات و باشگاه ورزشی به‌صورت دستی مدیریت می‌شود که منجر به تداخل رزروها و نارضایتی ساکنین می‌گردد.
\فقره \textbf{اطلاع‌رسانی ناکارآمد:} اطلاعیه‌های ساختمان از طریق برگه‌های چاپی یا گروه‌های پیام‌رسان غیررسمی منتشر می‌شوند که آرشیو منظمی ندارند.
\فقره \textbf{کاهش بهره‌وری مدیر ساختمان:} مدیر وقت زیادی را صرف پاسخ به تماس‌های تلفنی، ردیابی پرداخت‌ها و هماهنگی با کارکنان خدماتی می‌کند.
\فقره \textbf{عدم دسترسی آسان ساکنین به خدمات:} ساکنین برای هر درخواستی نیاز به مراجعه‌ی حضوری یا تماس با مدیر دارند.
\پایان{فقرات}

\قسمت{هدف پروژه}

هدف اصلی ساکنا ایجاد یک پلتفرم جامع تحت وب است که:

\شروع{شمارش}
\فقره مدیریت یکپارچه‌ی ساختمان را برای مدیران فراهم کند و کار دستی آن‌ها را به حداقل برساند.
\فقره تجربه‌ی دیجیتال روان برای ساکنین فراهم سازد تا بدون مراجعه‌ی حضوری امور خود را پیگیری کنند.
\فقره جریان کار کارآمد برای کارکنان خدماتی ایجاد کند تا وظایف ارجاع‌شده را شفاف ببینند.
\فقره شفافیت مالی و اطلاعاتی را در سطح مجتمع مسکونی بهبود دهد.
\فقره دسترسی ۲۴ ساعته‌ی ساکنین به خدمات و اطلاعات مجتمع را ممکن سازد.
\پایان{شمارش}

\قسمت{نیازسنجی}

بررسی وضعیت موجود نشان می‌دهد که:

\شروع{فقرات}
\فقره اکثر مجتمع‌های مسکونی ایران فاقد هرگونه ابزار دیجیتال برای مدیریت ساختمان هستند.
\فقره مدیران ساختمان به‌طور میانگین بخش قابل‌توجهی از وقت خود را صرف امور اداری می‌کنند که قابل اتوماسیون است.
\فقره ساکنین تمایل دارند امور روزمره‌ی ساختمانی را آنلاین انجام دهند، اما ابزار مناسبی در اختیار ندارند.
\فقره رشد نفوذ اینترنت و استفاده از گوشی‌های هوشمند در ایران، زمینه را برای پذیرش راه‌حل‌های دیجیتال آماده کرده است.
\پایان{فقرات}

\قسمت{فرصت محصول}

با توجه به گسترش آپارتمان‌نشینی در کلان‌شهرهای ایران و رشد تقاضا برای خدمات دیجیتال در مدیریت ساختمان،
ساکنا فرصت پاسخ‌گویی به نیاز صدها هزار مجتمع مسکونی را دارد.
راه‌حل‌های موجود در بازار یا ناقص هستند، یا پیچیدگی زیادی دارند، یا برای شرایط ایران طراحی نشده‌اند.
ساکنا با تمرکز بر سادگی، کارایی و یکپارچگی می‌تواند این خلأ را پر کند.
